\section{引言}\label{sec:introduction}
多变量时间序列记录了环境监测、生理传感、工业系统和金融市场中相互关联的动态过程。缺失观测以不同位置和持续时间打断这些过程，同时遮蔽局部动态与变量之间的联系。恢复缺失值需要融合这些互补维度上的可用信息。

现有方法通过循环估计、自注意力、扩散以及结构化时空表示建模不完整观测~\cite{brits,saits,csdi,imputeformer}。时间Patch为局部轨迹提供了有效的表示单元~\cite{patchtst}，变量注意力则捕捉不同测量过程之间的依赖~\cite{itransformer,t1}。这些进展引出了缺失修复的两个互补要求：在受损局部上下文之外寻找有用片段，以及将上下文表示转化为与可见测量一致的数值估计。

可用信息随缺失模式变化。同变量片段可以揭示时间变化，相关变量可以呈现共同转折，训练历史片段则可以提供重复出现的轨迹。不同证据在目标Patch内部的作用也随位置变化：靠近可见边界和位于缺口深处的点具有不同的局部支撑。因此，模型按来源组织Patch证据，并在单个缺失点的粒度上读取这些信息。

初始估计形成后，可见观测又为数值修正提供了依据。剩余测量锚定轨迹的水平与局部变化，独立的精恢复阶段可以结合初始估计与这些观测，纠正窗口中的剩余误差。由此形成粗到精设计：上下文重建提供初始轨迹，观测条件化精恢复进一步改善逐点数值。

基于上述考虑，本文提出证据引导的粗到精修复框架\Model{}，包含四个组件：（1）时间--变量上下文编码（TVCE）；（2）多来源Patch证据检索（MPER）；（3）逐点证据引导重建（PEGR）；（4）观测条件化残差精恢复（OCRR）。TVCE结合残差时间卷积与变量签名注意力；MPER将时间、跨变量和训练历史Patch组织为统一证据库；PEGR由各目标点分别读取证据，并结合Query与读取结果修正基础估计；OCRR进一步融合固定粗估计、可见观测及参考特征，生成精恢复窗口。

本工作的主要贡献如下：
\begin{itemize}
 \item 提出粗到精修复框架\Model{}，融合上下文重建、显式Patch证据和观测条件化修正，实现多变量时间序列缺失恢复。
 \item 设计多来源Patch证据检索与逐点重建机制，通过统一读取接口，为不同目标点提供互补的时间、跨变量及历史片段。
 \item 设计观测条件化残差精恢复器，将时间修正与可见观测参考结合。六个数据集的实验与配对组件分析评价整体框架及其恢复机制。
\end{itemize}

\begin{figure*}[t]
 \centering
 \IfFileExists{figures/evirefine_architecture.pdf}{\includegraphics[width=\textwidth]{figures/evirefine_architecture.pdf}}{%
 \fbox{\parbox{.94\textwidth}{\centering\vspace{4mm}
 \textbf{\Model{}架构}\\[2mm]
 TVCE＋MPER\\
 $\longrightarrow$PEGR＋置信度加权重叠融合\\
 $\longrightarrow$OCRR
 \vspace{4mm}}}}
 \caption{\Model{}结构概览。TVCE编码点级上下文，MPER构建Patch证据，PEGR生成粗重建，OCRR结合可见观测进行精恢复。}
 \label{fig:architecture}
\end{figure*}
